\section{Related Work}\label{sec:related}

\paragraph{Capabilities and attenuation.} Dennis and Van Horn describe a
computation's authority as a list of capabilities, where a process ``may grant
some of its capabilities to the inferior sphere'' but cannot pass one ``not
implied by a capability present in the higher sphere''~\cite{dennis}. Miller's
principle of least authority and attenuation by interposition~\cite{miller},
SPKI's intersection of authorization tags along a certificate
chain~\cite{spki} and Capsicum's rule that a derived capability's rights
``must be a subset of the original rights''~\cite{capsicum} are the same
monotonicity that Theorems~\ref{thm:atten} and~\ref{thm:chain} state for
grants. Machine-checked versions of it are older and deeper than ours: seL4's
integrity proof covers its capability-based access control and holds ``via
refinement for the C implementation of the kernel''~\cite{sel4Integrity}, and
CHERI's proofs establish ``a reachable-capability monotonicity property'' of
the instruction set~\cite{cheriProofs}. Our monotonicity theorems are
immediate in the model; what is new is enforcing one verified evaluator's
verdicts at an agent host and inside a smart account's authorization. The
relay is a deputy holding the custody key for many agents; it avoids the
confused-deputy problem~\cite{confusedDeputy} by checking the requesting
agent's grant, not its own authority.

\paragraph{Caveat tokens.} Macaroons embed caveats, ``predicates that restrict
the macaroon's authority'', and let any holder derive a narrower credential by
adding one, with chained HMACs giving ``the inability for adversaries to remove
caveats''~\cite{macaroons}. Biscuit expresses checks in Datalog and uses
public-key signatures, so a holder ``cannot remove existing blocks without
invalidating the signature''~\cite{biscuit}. UCAN requires that ``each direct
delegation MUST either directly restate or attenuate (diminish) its
capabilities''~\cite{ucan}, and its delegation policies constrain an
invocation's arguments while leaving stateful conditions ``out of scope of the
UCAN Policy language''~\cite{ucanDelegation}. ZCAP capabilities ``inherit the
restrictions from all caveat properties of their parents''~\cite{zcapld}.
These are bearer credentials whose attenuation is enforced cryptographically;
what a caveat means is left to the verifier's code. Grants are held by the
enforcer and used by a holder key, the evaluator that gives caveats their
meaning is proven to compute the model, and budgets add the stateful
condition these token languages leave out: spending charged up the delegation
tree.

\paragraph{Smart-account delegation and session keys.} Account
abstraction lets ``each account to supply its own validation
logic''~\cite{erc4337}, and modular accounts add validator and hook
modules~\cite{erc7579}. ERC-7710 standardizes redeeming a delegation through
a delegation manager, leaving ``the structure of the
\code{\_permissionContexts}'' to each manager and recommending audits; it
names granting ``bounded permissions to AI agents'' as a
motivation~\cite{erc7710}. ERC-7715 lets an application request permissions
from a wallet, which ``MAY attenuate(reduce or increase)'' them~\cite{erc7715}.
MetaMask's delegation framework enforces caveats with enforcer contracts whose
hooks run ``from leaf to root delegation'' on redemption; a re-delegation can
add caveats, and ``to understand what the delegation can do, it is necessary to
analyze all the enforcers being used in the entire delegation
chain''~\cite{metamaskDelegation}. Its enforcers include per-delegation and
periodic transfer limits, and it is assured by audits; we found no formal
verification. Smart Sessions attaches policies with argument conditions and
ERC-20 spending limits to session keys~\cite{smartSessions}, and Zodiac Roles
evaluates condition trees over calldata with refillable
allowances~\cite{zodiacRoles}. On Stellar, OpenZeppelin's smart accounts check
context rules with policy contracts whose \code{enforce} ``both validates
conditions and applies state changes''~\cite{ozStellarAccounts}, and an agent
wallet example caps an agent's spending per period on one
token~\cite{ozAgentWallet}. Google's Agent Payments Protocol binds an agent's
payment to the user's intent with signed mandates, whose constraints ``allow
the verifier to verify that the Checkout and Payment match the User's
intent''~\cite{ap2}. These systems express richer conditions than our closed
caveat set, and several have the periodic or refillable budgets we lack. In
each, a caveat's meaning is the code of its enforcer or policy contract,
assured by audit and analysed chain by chain; here one evaluator gives every
caveat its meaning at every boundary, and attenuation, chaining and budget
safety are theorems about that evaluator's translation. The grant signer also
relies on Soroban handing a delegated signer every context of the account's
authorization~\cite{cap7101}, which a policy attached to one rule does not
see.

\paragraph{Custody policy engines.} Dfns evaluates amount and recipient rules
on plain value transfers only~\cite{dfnsPolicies}. Other engines decode
contract calls. Fordefi's policy conditions can specify ``the method's
parameters and their values''~\cite{fordefiPolicies}, with byte-string, array
and tuple parameters since May 2026~\cite{fordefiChangelog}; Turnkey's policy
language reads ``the parsed Ethereum transaction payload'' and, once a
contract's interface is uploaded, its function calls and
arguments~\cite{turnkeyPolicies,turnkeyContracts}; and Fireblocks' policy rules
match contract methods~\cite{fireblocksPolicyApi}, and its own comparison
claims ``method-level allow and deny by ABI function'' and transaction
simulation~\cite{fireblocksPolicy}. The engines whose documentation we read
describe rules over EVM-style calls; we found none that evaluates a Soroban
authorization tree. Grants differ from all four in three ways: holders
delegate by attenuation, the evaluator is verified, and the same policy is
enforced again on chain by the account itself, so the custody key is not the
only check.

\paragraph{Verified policy evaluation and Rust verification.} Cedar's authors
proved properties of a Lean model of the language and ``used extensive
differential random testing to ensure that the implementation matches the
model''~\cite{cedar}; they describe formally verifying ``the equivalence of the
deployed code and model'' as the alternative they did not take~\cite{cedarHow}.
Zelkova encodes AWS access policies into SMT to analyse them~\cite{zelkova}.
Our policy language is far smaller than Cedar's, and in exchange the deployed
Rust's translation is proven equal to the model through Charon and
Aeneas~\cite{charon,aeneas}. Rust can also be verified in place: Verus with
SMT and specifications written in Rust~\cite{verus}, Creusot through
Why3~\cite{creusot}, Prusti through Viper~\cite{prusti}, and hax by
translating to several provers' languages~\cite{hax}. We chose a functional
translation to Lean for two reasons: the readable model and the translation
live in one language, so the equivalence is an ordinary Lean theorem, and the
proofs use Lean's library and kernel rather than an SMT solver's.

\paragraph{WebAssembly sandboxing.} WebAssembly was designed with a formal
semantics~\cite{wasm}, mechanized in Isabelle~\cite{wasmcert}. vWasm and rWasm
compile Wasm into provably sandboxed code~\cite{bosamiya}, RLBox sandboxes
libraries inside Firefox with Wasm~\cite{rlbox}, and WaVe is a verified WASI
runtime ensuring that each file a sandbox accesses lies within its root
directory and that sockets reach only allow-listed addresses~\cite{wave}.
These isolate a module from resources it was not given. The agent host instead
gives an agent's tool calls a per-call check against a delegable grant, and
runs the checking evaluator itself as a component that imports
nothing~\cite{componentModelDocs}; it trusts Wasmtime~\cite{wasmtime} rather
than verifying the runtime, and confines no agent.

\paragraph{Agent security.} CaMeL tracks data flow through an agent's program
and enforces capability policies, written as Python functions, when tools are
called~\cite{camel}. FIDES ``propagates labels in messages, actions, tool calls
and results'' and executes consequential actions only when a policy over those
labels allows~\cite{fides}; both address the information flow that grants
leave out. Progent expresses privilege as ``symbolic rules over tool names and
arguments'' in JSON Schema and uses an SMT solver to decide whether a policy
update narrows or expands the agent's actions~\cite{progent}, and AgentSpec
lets users write rules with ``triggers, predicates, and enforcement
mechanisms'' evaluated at runtime~\cite{agentspec}; CaMeL and Progent evaluate
on AgentDojo~\cite{agentdojo}. A theory of what such gates can enforce shows
that ``deterministic pre-execution gates enforce exactly the nonempty
register-representable safety policies''~\cite{whatCanBeEnforced}; a grant
check is such a gate. nono
confines agent processes at the kernel with path, process and network rules
down to the HTTP method and path~\cite{nono}. Grants share Progent's monotonic
narrowing and nono's placement in front of the agent's resources, but take
their checks down to argument values, carry them past the agent's process into
the smart account that executes the transaction, and prove the evaluator that
decides them.
